\bookchapter{Capítulo 34: Guía detallada de conectividad móvil-automóvil — La fusión del teléfono y el coche}

\chapterintro{Usa el teléfono como nivel de burbuja}

\booksection{: La app de Brújula del iPhone tiene una función de nivel (desliza hacia la izquierda). Android también tiene apps de nivel; útil para colgar cuadros o estanterías.}

Usa el flash del teléfono como linterna

\begin{itemize}
  \item \textbf{: Activa la linterna en el Centro de control (iPhone) o en la barra de notificaciones (Android). Útil en apagones o al buscar algo en la oscuridad.}Usa la alarma del teléfono como recordatorio
  \item \textbf{: Además de despertarte, la alarma puede servirte de recordatorio para tomar medicina, recoger a los niños o una reunión importante.}Usa el cronómetro y el temporizador para la gestión del tiempo
  \item \textbf{: Usa el cronómetro para medir el tiempo de tareas (técnica Pomodoro) y el temporizador para limitar el tiempo de uso del teléfono o de la cocina.}Usa la grabadora de voz para tomar notas rápidas
\end{itemize}

\booksection{: Si no puedes escribir, graba tus ideas en voz. Útil para reuniones, ideas repentinas o cuando conduces.}

Usa la app de Tiempo para planificar el día

\begin{center}
\begin{tabularx}{\linewidth}{|X|X|X|X|X|}
\hline
\textbf{: Mira el tiempo por la mañana para decidir la ropa y el paraguas. Activa las alertas de lluvia para no pillarte fuera.} & \textbf{Usa el teléfono como cuentakilómetros} & \textbf{: Las apps de salud (Salud en iPhone, Google Fit en Android) registran tus pasos diarios. Útil para controlar la actividad física.} & \textbf{Usa "Buscar mi teléfono" en caso de pérdida} & \textbf{: Si pierdes el teléfono, usa "Buscar mi iPhone" o "Buscar mi dispositivo" desde otro dispositivo para localizarlo, bloquearlo o borrar los datos.} \tabularnewline
\hline
Apple CarPlay & Usa el modo avión para cargar más rápido & : Al cargar el teléfono, activa el modo avión para que cargue más rápido (al no buscar red ni Wi-Fi, consume menos mientras carga). & Usa el teléfono como disco duro portátil & : Con un cable USB-C/Lightning, puedes transferir archivos entre el teléfono y el ordenador. Útil para llevar documentos o fotos grandes. \tabularnewline
Android Auto & VI. Consejos para usuarios mayores & El teléfono puede ser una gran herramienta para los mayores, pero requiere algunos ajustes para ser más accesible: & Activa el "modo sencillo" & : Muchos teléfonos Android tienen un "modo sencillo" que agranda los iconos y simplifica la interfaz. iPhone no tiene esta función, pero se puede lograr un efecto similar ajustando el tamaño de los iconos y el texto. \tabularnewline
Huawei HiCar & Aumenta el tamaño del texto & : Ajustes > Pantalla > Tamaño del texto. Aumenta el tamaño del texto para que sea más fácil de leer. iPhone también tiene "Texto en negrita" que mejora la legibilidad. & Activa la lectura en voz alta & : El iPhone tiene una función de "Lectura en voz alta" (Ajustes > Accesibilidad > Contenido leído). Android también la tiene; permite que el teléfono lea en voz alta el texto de la pantalla. \tabularnewline
Baidu CarLife & Activa el "modo de emergencia" y los contactos de emergencia & : Configura contactos de emergencia en el teléfono (Ajustes > Emergencia). En caso de necesidad, cualquiera puede llamar a tus contactos de emergencia desde la pantalla de bloqueo. & Desactiva las notificaciones no esenciales & : El iPhone tiene una función de "Lectura en voz alta" (Ajustes > Accesibilidad > Contenido leído). Android también la tiene; permite que el teléfono lea en voz alta el texto de la pantalla. \tabularnewline
\hline
\end{tabularx}
\end{center}

\booksection{: Los mayores pueden confundirse con demasiadas notificaciones. Desactiva las notificaciones de apps que no sean esenciales (noticias, juegos, compras).}

VII. Diez habilidades para usuarios avanzados

\begin{itemize}
  \item \textbf{Estas diez habilidades te llevan de "usar el teléfono" a "dominar el teléfono":}Aprende a usar "accesos directos" (Shortcuts)
  \item \textbf{: Los "Atajos" del iPhone y los "Atajos" de Android permiten automatizar tareas complejas. Por ejemplo, un atajo que, al salir del trabajo, envía un mensaje a tu pareja, activa la navegación y sube el volumen.}Usa "Automaciones" para tareas programadas
  \item \textbf{: Las automatizaciones del iPhone pueden ejecutar acciones según la hora, la ubicación o un evento. Por ejemplo, al llegar a casa, conecta el Wi-Fi y reproduce música; al salir del trabajo, activa el modo no molestar.}Domina la "pantalla dividida" y la "ventana flotante"
\end{itemize}

\booksection{: En Android, usa la pantalla dividida para hacer dos cosas a la vez; usa la ventana flotante para ver un vídeo mientras usas otras apps. En iPhone, usa “Image in Picture” para el mismo propósito.}

Usa “texto escaneado” para digitalizar documentos

\begin{itemize}
  \item \textbf{: El iPhone puede extraer texto de cualquier imagen o foto (función "Texto en vivo"). Android también lo soporta (Google Lens o funciones nativas). Útil para digitalizar tarjetas de visita, carteles y documentos.}Aprende a usar el "modo escritorio"
  \item \textbf{: Algunos teléfonos Android (Samsung DeX, Huawei) tienen un "modo escritorio" que, al conectar una pantalla, convierte el teléfono en un ordenador. Útil para oficina ligera y presentaciones.}Usa "Buscar mi teléfono" para localizarlo
  \item \textbf{: Si pierdes el teléfono, usa "Buscar mi iPhone" o "Buscar mi dispositivo" desde otro dispositivo para localizarlo, bloquearlo o borrar los datos de forma remota.}Domina la "grabación de pantalla"
  \item \textbf{: El iPhone y Android tienen grabación de pantalla integrada. Útil para grabar tutoriales, guardar vídeos que no se pueden descargar o reportar bugs.}Usa "modo de una mano" para teléfonos grandes
\end{itemize}

\booksection{: Si el teléfono es demasiado grande, activa el modo de una mano para reducir la pantalla y operar con un solo dedo.}

Aprende a usar “transferencia de archivos rápida”

\begin{itemize}
  \item \textbf{: AirDrop del iPhone permite transferir archivos rápidamente entre dispositivos Apple. Android tiene "Quick Share" o "Huawei Share" para lo mismo. Útil para compartir fotos y vídeos grandes.}Usa el teléfono como mando a distancia universal
  \item \textbf{: Muchos teléfonos tienen emisor infrarrojo y pueden usarse como mando de la TV, el aire acondicionado, etc. Descarga la app "Mando" o similar.}VIII. Diez consejos para la batería
  \item \textbf{Recomendacion}La batería es la parte del teléfono que más se degrada con el uso. Estos diez consejos prolongan su vida útil:
\end{itemize}

\begin{tcbwarn}{Alerta de seguridad}
No dejes que la batería llegue a 0\%

\end{tcbwarn}

\booksection{: Las baterías de litio no les gusta llegar a 0\% ni quedarse al 100\% mucho tiempo. Intenta cargar el teléfono cuando llegue al 20-30\% y desenchufarlo al 80-90\%.}

\begin{itemize}
  \item \textbf{No uses el teléfono mientras se carga}: Usar el teléfono mientras se carga genera calor extra y acelera el envejecimiento de la batería. Si tienes que usarlo, al menos no juegues ni grabes vídeos.
  \item \textbf{Evita las temperaturas extremas}: El calor es el mayor enemigo de la batería. No dejes el teléfono al sol ni en un coche caliente en verano. El frío extremo también daña la batería, aunque de forma temporal.
  \item \textbf{Activa la "carga optimizada"}: El iPhone tiene "Carga optimizada" que aprende tus hábitos de carga y retrasa la carga al 80\% hasta justo antes de que lo necesites. Android también tiene funciones similares.
\end{itemize}

\begin{tcbinfo}{No cargues el teléfono al 100\% todas las noches}
: Mantener la batería al 100\% durante mucho tiempo la estresa. Si puedes, cárgala al 80-90\% y desenchúfala.

\end{tcbinfo}

\begin{tcbpractice}{Desactiva el GPS, Bluetooth y Wi-Fi cuando no los uses}
\begin{enumerate}
  \item : Estos servicios consumen batería. Desactívalos cuando no los necesites; en Android, usa el panel de ajustes rápidos para hacerlo de un toque.
  \item Reduce el brillo de la pantalla
  \item : La pantalla es el mayor consumidor de batería. Activa el brillo automático o ajústalo manualmente al nivel mínimo cómodo.
  \item Desactiva la actualización en segundo plano de las apps
\end{enumerate}

\end{tcbpractice}

\begin{tcbsummary}{Resumen del capítulo}
\begin{itemize}
  \item : Ajustes > General > Actualización en segundo plano (iPhone) o Ajustes > Apps (Android). Desactiva la actualización en segundo plano de las apps que no necesiten estar siempre actualizadas.
  \item Usa el modo avión en zonas con mala cobertura
  \item : Si estás en una zona con mala cobertura, el teléfono busca constantemente la red y consume mucha batería. Activa el modo avión o desactiva los datos móviles para ahorrar batería.
  \item Cambia la batería cuando su salud baje del 80\%
  \item : El iPhone muestra la salud de la batería (Ajustes > Batería > Salud de la batería). Si baja del 80\%, considera cambiarla en un servicio técnico oficial.
  \item IX. Diez consejos para la cámara
\end{itemize}

\end{tcbsummary}
